\subsection{Planejamento da trajetória e órbita de fase}

A estratégia de transferência inicial utiliza a transferência de Hohmann para transição entre as órbitas circulares e coplanares, permitindo determinar os instantes da manobra impulsiva, janelas de lançamento e o $\Delta v$ requerido para a inserção na órbita de faseamento. 
A órbita de faseamento serve para ajustar a defasagem angular entre o \emph{chaser} e o \emph{target}, possibilitando que a aproximação se inicie com alinhamento de fase adequado.

O procedimento segue três etapas principais:
\begin{enumerate}[label=\roman*.]
    \item \textbf{Transferência de Hohmann}: cálculo dos parâmetros da manobra impulsiva e do tempo de voo até a nova órbita circular;
    \item \textbf{Órbita de faseamento}: determina a defasagem angular e do período orbital necessário para sincronização com o \emph{target};
    \item \textbf{Janela de aproximação}: definição das condições de início da fase de \emph{close approach}, considerando o alcance máximo e o tempo relativo de chegada.
\end{enumerate}

Os erros de sincronização ou os atrasos na execução das manobras impulsivas resultam em desvios de fase indesejados e no maior consumo de propelente para correção.  
Durante a transição da órbita de fase para a aproximação de curta distância, são utilizados os modelos lineares de Hill/Clohessy--Wiltshire para o movimento relativo.
